\documentclass[11pt]{article}
\usepackage[margin=2.2cm]{geometry}
\usepackage{booktabs,longtable,array,ragged2e}
\usepackage{amsmath}
\usepackage[colorlinks=true,linkcolor=blue,urlcolor=blue]{hyperref}
\usepackage{xcolor}
\setlength{\parskip}{4pt}
\setlength{\parindent}{0pt}
\renewcommand{\arraystretch}{1.15}

\title{\bfseries Supporting Information\\[4pt]
\large The p53 pulse encoder counts by topological necessity:\\
theory, evidence, and the molecular origin of the 5.5-hour clock}
\author{Huan Li\\ \small Shanghai Liangta Biotech Co., Ltd., Shanghai 201899, P.\ R.\ China}
\date{October 2026}

\begin{document}
\maketitle

\noindent This Supporting Information documents the computational output of
the accompanying repository (\texttt{p53-pulse-encoder}, release v1.0.2): the
analysis scripts, the archived result file for each analysis that produces
one, and the provenance of third-party data. Numbers below are transcribed
from the archived result files (\texttt{results/}); where a script writes only
to its console, the pipeline map says so. Scripts are plain Python~3
(numpy/scipy/matplotlib). The digitised Batchelor Fig.~1G/H extraction is shipped in
\texttt{data\_external/batchelor2011/} and is the direct input of code52; its
values and error bands are also reported in the main text.

\tableofcontents
\newpage

%% ----------------------------------------------------------------
\section{Pipeline map: scripts and outputs}

\begin{longtable}{p{4.6cm} p{6.6cm} p{4.6cm}}
\toprule
Script & Computes & Archived output \\
\midrule
\endhead
\texttt{code13\_\allowbreak{}p53\_\allowbreak{}fisher\_\allowbreak{} stratification\_\allowbreak{}S7} & Fisher stratification of the pulse generator, S7 slice & \texttt{p53\_\allowbreak{}channel\_\allowbreak{}census.csv} \\
\texttt{code50\_\allowbreak{}p53\_\allowbreak{}cascade\_\allowbreak{} framework\_\allowbreak{}selfcheck} & internal consistency self-check of the cascade framework & console log \\
\texttt{code51r1--r4\_\allowbreak{}p53\_\allowbreak{}fisher\_\allowbreak{} audit} & Fisher identifiability audit of the phenomenological pulse model; r4 adds noise calibration and an $M{=}4000$ central-difference control & code51r2/r3 results JSONs (\texttt{code51r2\_\allowbreak{}p53\_\allowbreak{}Fisher\_\allowbreak{}audit\_\allowbreak{}results.json}, \texttt{code51r3\_\allowbreak{}...json}), \texttt{code51r4A\_\allowbreak{}...json}, \texttt{code51r4B\_\allowbreak{}Fisher\_\allowbreak{}M4000.json}, three verdict figures (PNG) \\
\texttt{code52\_\allowbreak{}TrackC2\_\allowbreak{} discrimination\_\allowbreak{}diagram} & discrimination map between counting and amplitude encoders & \texttt{code52\_\allowbreak{}TrackC2\_\allowbreak{}results.json}, discrimination-map PNG \\
\texttt{code81\_\allowbreak{}p53\_\allowbreak{}full\_\allowbreak{}cascade\_\allowbreak{} audit} & full-cascade simulation (damage $\to$ fate), information identity, amplitude-share sweep & \texttt{code81\_\allowbreak{}cascade\_\allowbreak{}audit\_\allowbreak{}results.json}, five-panel PNG \\
\texttt{code82\_\allowbreak{}p53\_\allowbreak{}mechanism\_\allowbreak{} model\_\allowbreak{}audit} & mechanistic 6-species model; baseline reproduction, Fisher audit, mechanism tests & \texttt{code82\_\allowbreak{}mechanism\_\allowbreak{}audit\_\allowbreak{}results.json}, two four-panel PNGs \\
\texttt{code83\_\allowbreak{}why\_\allowbreak{}digital\_\allowbreak{} variational\_\allowbreak{}competition} & variational argument for digital pulses & console/tables \\
\texttt{code84\_\allowbreak{}discovery\_\allowbreak{} cascade\_\allowbreak{}model} & discovery-grade cascade variant & console/tables \\
\texttt{code85\_\allowbreak{}heterogeneity\_\allowbreak{} recheck} & heterogeneity robustness re-check & console/tables \\
\texttt{code86\_\allowbreak{}foci\_\allowbreak{}persistence\_\allowbreak{} trigger} & DSB-foci persistence as trigger & console/tables \\
\texttt{code87\_\allowbreak{}clock\_\allowbreak{}dose\_\allowbreak{} independent\_\allowbreak{}architecture} & dose-independence of the clock architecture & console/tables \\
\texttt{code88\_\allowbreak{}foci\_\allowbreak{}termination\_\allowbreak{} channel} & foci-termination channel analysis & console/tables \\
\texttt{code89\_\allowbreak{}p53\_\allowbreak{}molecular\_\allowbreak{} constant\_\allowbreak{}decomposition} & decomposition of the clock period into molecular phase constants & \texttt{molecular\_\allowbreak{}layer/code89\_\allowbreak{}...json} \\
\texttt{code90\_\allowbreak{}p53\_\allowbreak{}full\_\allowbreak{}cascade\_\allowbreak{} consistency\_\allowbreak{}audit} & 14 static link closures; onset-time decomposition & \texttt{molecular\_\allowbreak{}layer/code90\_\allowbreak{}...json} \\
\texttt{code91\_\allowbreak{}K2\_\allowbreak{}pulse\_\allowbreak{}width\_\allowbreak{} decomposition} & pulse-width (FWHM) molecular attribution & \texttt{molecular\_\allowbreak{}layer/code91\_\allowbreak{}...json} \\
\texttt{code92\_\allowbreak{}belgacem\_\allowbreak{}closedform\_\allowbreak{} crossvalidation} & reproduction of the closed-form period of arXiv:2605.23722 and reconciliation & \texttt{molecular\_\allowbreak{}layer/code92\_\allowbreak{}...json} \\
\texttt{code93\_\allowbreak{}tcs\_\allowbreak{}logistic\_\allowbreak{} correspondence} & event-time invariance under logistic time warp & \texttt{molecular\_\allowbreak{}layer/code93\_\allowbreak{}...json} \\
\texttt{cellline1\_bifurcation}\allowbreak\texttt{\_suite/} (scripts 1--3) & stochastic population simulation, bifurcation analysis, discrimination diagram (cell-line 1 suite cited in Methods) & figures \\
\multicolumn{3}{l}{\emph{Manuscript figures (all in \texttt{manuscript/})}}\\
\texttt{make\_\allowbreak{}fig1.py} & Fig.~1 encoder topology & \texttt{fig1\_\allowbreak{}topology.png} \\
\texttt{make\_\allowbreak{}fig2.py} & Fig.~2 channel census & \texttt{fig2\_\allowbreak{}census.png} \\
\texttt{make\_\allowbreak{}fig3.py} & Fig.~3 molecular decomposition (figure layer of code89; reads \texttt{code89\_\allowbreak{}...json}) & \texttt{fig3\_\allowbreak{}molecular.png} \\
\texttt{make\_\allowbreak{}fig4.py} & Fig.~4 closed-form cross-validation (figure layer of code92; reads \texttt{code92\_\allowbreak{}...json}) & \texttt{fig4\_\allowbreak{}closedform.png} \\
\bottomrule
\end{longtable}

Run metadata: code81/82 seed \texttt{20260923}; code89--93 seed
\texttt{20260925}; Fisher audits use relative step $h\in\{2\%,5\%\}$ with a
central-difference control at $M{=}4000$ trajectories.

%% ----------------------------------------------------------------
\section{Fisher identifiability audit (code51r1--r4)}
\label{sec:fisher}

Operating point $\theta_0$ of the phenomenological model
(Table~\ref{tab:theta0}); observables: $t_1$, $T$, $W$, amplitude, and the
pulse-count distribution.

\begin{table}[h]
\centering\small
\caption{Operating point used in all Fisher audits.}
\label{tab:theta0}
\begin{tabular}{lccccccc}
\toprule
$a$ & $b$ & $\epsilon$ & $s$ & $\tau_r$ (min) & $I_0$ & $\sigma$ \\
\midrule
0.7 & 0.8 & 0.08 & 20.0 & 20.0 & 0.6 & 0.06 \\
\bottomrule
\end{tabular}
\end{table}

\begin{table}[h]
\centering\small
\caption{Condition numbers of the Fisher matrix across audits.}
\label{tab:cond}
\begin{tabular}{lll}
\toprule
Audit & Setting & $\kappa(F)$ \\
\midrule
code51r1 & $h{=}2\%$ / $5\%$ & $\infty$ (exact null direction, $k_{\rm obs}$) \\
code51r2 & $h{=}2\%$ & $7.52\times10^{17}$ \\
code51r2 & $h{=}5\%$ & $1.38\times10^{4}$ \\
code51r3 & $h{=}2\%$ & $4.87\times10^{17}$ \\
code51r3 & $h{=}5\%$ & $4.37\times10^{6}$ \\
code51r4B & $M{=}4000$, central, $h{=}2\%$ & $2.13\times10^{18}$ \\
\bottomrule
\end{tabular}
\end{table}

Supporting quantities: mean pulse count $\bar N = 9.04$ (r1/r3), $8.33$
(r2), $9.01$ (r4B); $\mathrm{var}(t_1)=1.29\times10^{-4}$,
$\mathrm{var}(T)=1.09\times10^{-3}$, $\mathrm{var}(t_1/T)=0.119$.
Noise calibration (code51r4A): adopted $\sigma^\star=0.25$, giving
$\mathrm{CV}_T=0.111$, IPI $\mathrm{CV}=0.316$, valid fraction $1.0$,
$\bar N=10.39$. Channel capacity across 9 doses:
$I(D;N)=3.07$~bits vs $H(D)=3.17$~bits (code51r3).

\section{Counting-vs-amplitude discrimination (code52)}

Slope estimators on the Track-C2 discrimination diagram
(Table~\ref{tab:trackc2}).

\begin{table}[h]
\centering\small
\caption{Discrimination slopes (mean $\pm$ s.d., 95\% CI).}
\label{tab:trackc2}
\begin{tabular}{lccc}
\toprule
Arm & mean & s.d. & 95\% CI \\
\midrule
$\alpha_{\rm NCS}$ & $-0.0008$ & 0.113 & $(-0.225,\ 0.222)$ \\
$\alpha_{\rm UV}$ & $0.981$ & 0.114 & $(0.777,\ 1.224)$ \\
\bottomrule
\end{tabular}
\end{table}

The $\tau_r$ inversion point sits at $7.7$~h.

\section{Full-cascade audit (code81)}
\label{sec:code81}

Simulation: 7 doses $\times$ 4000 cells/dose. The 20 frozen constants are
listed in Table~\ref{tab:c81const}.

\begin{table}[h]
\centering\small
\caption{Frozen constants of the code81 cascade (verbatim from the archive).}
\label{tab:c81const}
\begin{tabular}{llll}
\toprule
Constant & Value & Constant & Value \\
\midrule
$K0_{\rm DSB}$ & 35.0 & $Q_{21}$ / $\tau_{21}$ & 100.0 / 10.0 \\
$C_{\rm CLUSTER}$ & 50.0 & $Q_{\rm PU}$ / $\tau_{\rm PU}$ & 100.0 / 4.0 \\
$N_{\rm MAX}$ & 12 & PUMA $n_{\rm thr}$ & 3 \\
$T_{\rm PERIOD}$ (h) & 5.5 & $\theta_{\rm PUMA}$ & 800.0 \\
$\mu_D/T$ & 1.0 & $R0_{\rm FATE}$ & 0.15 \\
JIT$_{\rm IPI}$ & 0.085 & $\theta_{\rm P21}$ & 1500.0 \\
$A_0$ & 1.0 & $B_{\rm GAIN}$ & 0.3 \\
$A_{\rm CV}$ & 0.002 & $\eta_{\rm OBS}$ & 0.02 \\
$W$ (h) & 48.0 & $\alpha_E$ & 7.0 \\
\bottomrule
\end{tabular}
\end{table}

\begin{table}[h]
\centering\small
\caption{Fate table (fraction of 4000 cells) vs dose.}
\label{tab:fate}
\begin{tabular}{lccc}
\toprule
Dose (Gy) & survival & arrest & apoptosis \\
\midrule
0.1 & 1.0000 & 0.0000 & 0.0000 \\
0.25 & 1.0000 & 0.0000 & 0.0000 \\
0.5 & 1.0000 & 0.0000 & 0.0000 \\
1.0 & 0.9948 & 0.0053 & 0.0000 \\
2.5 & 0.0047 & 0.9952 & 0.0000 \\
5.0 & 0.0015 & 0.9688 & 0.0297 \\
10.0 & 0.0020 & 0.0088 & 0.9892 \\
\bottomrule
\end{tabular}
\end{table}

Information identity: $H(D)=2.8074$~bits decomposes into the five layer
contributions with closure residual $0.0$; the weakest layer is L1
(sensing), losing $0.872$~bits. Amplitude share of dose information at
kinetic-noise levels $\sigma_k = 0.0/0.3/0.6$: $34.5\% / 14.8\% / 9.9\%$
of the counting information ($I_{\rm count}=1.631$~bits throughout).

\section{Mechanistic-model audit (code82)}

Six-species model (ATM$^*$, p53, mdm2 mRNA, Mdm2, wip1 mRNA, Wip1);
23 parameters, audited in the 2025 and 2017 parameter sets
(Table~\ref{tab:params82}).

\begin{table}[h]
\centering\footnotesize
\caption{Parameter sets (code82 archive).}
\label{tab:params82}
\begin{tabular}{lcccccccccccc}
\toprule
Set & $A$ & $P$ & $C$ & $g$ & $d_{\rm AM}$ & $T_m$ & $T_M$ & $T_w$ & $T_W$ & $d_A$ & $d_P$ & $d_m$ \\
\midrule
2025 & 30.5 & 22 & 1.4 & 2.5 & 20 & 1.2 & 4 & 1.2 & 1 & 0.16 & 0.1 & 1 \\
2017 & 30.5 & 22 & 1.4 & 2.5 & 20 & 1.0 & 4 & 1.0 & 1 & 0.16 & 0.1 & 1 \\
\midrule
 & $d_M$ & $d_w$ & $d_W$ & $k_A$ & $k_{\rm WA}$ & $k_{\rm MP}$ & $k_{\rm Pm}$ & $k_{\rm Pw}$ & $R$ & $S_{\max}$ & $\gamma$ & \\
\cmidrule{2-12}
2025 & 2 & 1.3 & 2.3 & 0.5 & 0.14 & 0.15 & 1 & 1 & 2 & 0.2 & 9 & \\
2017 & 2 & 1.3 & 2.3 & 0.5 & 0.14 & 0.15 & 1 & 1 & 2 & 0.2 & 9 & \\
\bottomrule
\end{tabular}
\end{table}

Ground state (both sets): ATM$^*$ 0.0353, p53 0.2180, mdm2 0.2147,
Mdm2 0.3174, wip1 0.1652, Wip1 0.0718.

\begin{table}[h]
\centering\small
\caption{R1 baseline reproduction: pulse count, period, amplitude.}
\label{tab:r1}
\begin{tabular}{llccc}
\toprule
Set & Damage input & pulses & period (h) & amplitude \\
\midrule
2025 & log & 17 & 5.484 & 2.783 \\
2025 & saturating & 17 & 5.753 & 2.788 \\
2017 & log & 14 & 6.831 & 3.408 \\
2017 & saturating & 14 & 6.681 & 3.407 \\
\bottomrule
\end{tabular}
\end{table}

Damage-signal integral $S=0.4988$ (log) / $0.1982$ (saturating).
The stiffest Fisher axis is the Wip1-production axis
($T_w\,0.27 + T_W\,0.26 + P\,0.11 + d_W\,0.10$).

\section{Molecular decomposition of the 5.5 h clock (code89)}

Angular frequency $\omega = 1.1424$~rad/h.

\begin{table}[h]
\centering\small
\caption{Phase-constant decomposition vs measurement.}
\label{tab:code89}
\begin{tabular}{lcc}
\toprule
Quantity & Predicted & Target \\
\midrule
$T_1$: K3 mdm2 lag (h) & $1.70$ & $2.0 \pm 0.5$ \\
$T_2$: K4 wip1 delay (h) & $1.13$ & $1.25$ \\
$T_3$: K1 period (h) & $5.95$ & $5.48$ \\
\bottomrule
\end{tabular}
\end{table}

Per-species phase shares of the period: p53 $30.4\%$, ATM$^*$ $24.3\%$,
mdm2 mRNA $18.2\%$, wip1 mRNA $15.3\%$, Wip1 $9.6\%$, Mdm2 $2.2\%$
(sum of phases $= 4.474$~rad). Molecular closure register: $T_W$
(translation, universal $\sim$4 aa/s) and $d_W$ (Wip1 degradation,
0.7--2.3/h across sources) are measurable constants; $T_w$ and $P$ are
registered gaps.

\section{Full-cascade consistency (code90)}

14 static links audited: \textbf{7 closed, 3 tension, 4 gap.}
Onset decomposition at the step signal ($S=0.4988$, basal $0.0401$):
sensing contributes only $0.83$~h in total ($\gamma$H2AX segment 0.50~h,
ATM half-rise 0.19~h, p53 threshold crossing 0.14~h); the first p53 peak
arrives at $3.21$~h, i.e.\ the bulk of the measured $t_1$ (MCF7 band
2--3~h) is the accumulation ramp ($2.38$~h), not sensing delay.

\section{Pulse-width decomposition (code91)}

At DSB $=100$: measured FWHM $W=2.779$~h decomposes as
p53 half-rise $\to$ ATM$^*$ extinction ($0.79$~h) $+$ collapse tail
($1.99$~h). The plateau segment ($1.54$~h) is set by Wip1 accumulation;
the collapse tail by Mdm2 recovery kinetics ($d_M/d_{\rm AM}$) --- the
dominant molecular attribution of width is Mdm2 recovery, not the Wip1
threshold crossing itself. Dose scan: width $\mathrm{CV}=0.055$.

\section{Closed-form cross-validation (code92)}

Reproduction of the closed form of Belgacem
(arXiv:2605.23722, Section~9): $AB = 1.7204$ inverted from the manuscript
constants vs $1.72$ published; $\tau_c = 1.000$~h; closed-form period
$T=5.64$~h vs $5.6$~h published and $5.5$~h measured.
Reconciliation row: measured $5.5$~h; Belgacem closed form $5.64$~h;
Monke simulation $5.48$~h; Monke Jacobian $5.95$~h; our $\tau_2$
closed-form evaluation $8.40$~h (registered as a tension, see the main
text).

\section{TCS--logistic correspondence (code93)}

Event-time invariance under the logistic time warp: 7 pulses before and
after warping, maximal event-time shift $0.0$~h, period unchanged at
$5.4802$~h, while the amplitude CV inflates by $+42.3\%$
($0.308 \to 0.438$).

\section{Third-party data provenance}

No large third-party archive is re-hosted. Extracted tables in
\texttt{results/} trace to: Batchelor et al., \textit{Mol.\ Syst.\ Biol.}
7, 488 (2011); Stewart-Ornstein \& Lahav, \textit{Sci.\ Signal.} 10,
eaah6671 (2017) (PMC5504473); Geva-Zatorsky et al., \textit{Mol.\ Syst.\
Biol.} 2, 33 (2006); Loewer et al., \textit{Cell} 142, 89--100 (2010)
(PMC3003696); I.~Belgacem, arXiv:2605.23722 (2026). Molecular-layer
constants (delays, half-lives) trace to Wang et al., \textit{IET Syst.\
Biol.} 13, 180 (2019); Monk, \textit{Curr.\ Biol.} 13, 1409 (2003);
Batchelor et al., \textit{Mol.\ Cell} 30, 277 (2008); M\"onke et al.,
\textit{Sci.\ Rep.} 7, 46571 (2017). Full
access notes in \texttt{data\_external/PROVENANCE.md}.

\end{document}
